% =====================================================================
%  Anatomía de un LLM — Fase I — Módulo 1
%  Modelos de lenguaje probabilísticos
%  Compilar: latexmk -xelatex M01-modelos-de-lenguaje.tex
% =====================================================================
\documentclass[11pt,a4paper,twoside,openany]{book}
\usepackage{../estilo-apuntes}

\newcommand{\ApuntesCurso}{Anatomía de un LLM}
\newcommand{\ApuntesModulo}{Módulo 1}

\hypersetup{pdftitle={Módulo 1. Modelos de lenguaje probabilísticos},pdfauthor={Anatomía de un LLM}}

\begin{document}
\frontmatter
% ------------------------------------------------------------------ Portada
\begin{titlepage}
\centering
\vspace*{2.5cm}
{\large\scshape Anatomía de un LLM}\par
\vspace{0.4cm}
{\normalsize Fundamentos matemáticos de los modelos de lenguaje}\par
\vspace{2.2cm}
\rule{\textwidth}{0.8pt}\par\vspace{0.5cm}
{\small\scshape Fase I \textperiodcentered\ Fundamentos matemáticos}\par\vspace{0.3cm}
{\Huge\bfseries Módulo 1\par}\vspace{0.3cm}
{\LARGE Modelos de lenguaje probabilísticos\par}
\vspace{0.5cm}\rule{\textwidth}{0.4pt}\par
\vspace{1.6cm}
{\large Apuntes del curso}\par
\vfill
{\small Edición de octubre de 2026}\par
\end{titlepage}

% ------------------------------------------------------- Presentación
\chapter*{Presentación del módulo}
\addcontentsline{toc}{chapter}{Presentación del módulo}
\markboth{Presentación del módulo}{}

Un modelo de lenguaje es una distribución de probabilidad sobre secuencias de símbolos. Los grandes modelos de lenguaje actuales (Claude, Gemini, Grok y otros) son modelos de este tipo: dado un contexto, asignan una probabilidad a cada posible símbolo siguiente. Este módulo establece el marco probabilístico común a todos ellos y lo desarrolla en el caso más sencillo, el modelo de $n$-gramas, en el que cada cantidad puede calcularse a mano.

\section*{Objetivos}
Al finalizar el módulo, el estudiante debe ser capaz de:
\begin{enumerate}[label=\arabic*.]
  \item definir con precisión vocabulario, tokenización y modelo de lenguaje;
  \item enunciar y demostrar la regla de la cadena y formular la hipótesis de Markov;
  \item obtener el estimador de máxima verosimilitud de un modelo de bigramas y aplicar técnicas de suavizado;
  \item definir la función softmax con temperatura, demostrar sus propiedades y describir los métodos de muestreo;
  \item calcular la entropía cruzada y la perplejidad de un modelo y justificar su uso como criterio de entrenamiento y evaluación.
\end{enumerate}

\section*{Requisitos previos}
Probabilidad discreta elemental (espacio muestral, sucesos, probabilidad condicionada), logaritmos y exponenciales, sumatorios y productorios, y derivadas parciales. Los multiplicadores de Lagrange se usan en los temas~3 y~5; el apéndice de soluciones incluye todos los pasos.

\section*{Organización}
\begin{center}
\begin{tabular}{@{}cll@{}}
\toprule
Tema & Contenido & Fuentes principales \\
\midrule
1 & Vocabulario, tokenización y modelos de lenguaje & \cite{jm3}, \cite{cs224n} \\
2 & Regla de la cadena y modelos de Markov & \cite{bt}, \cite{mit6041}, \cite{jm3} \\
3 & Máxima verosimilitud y suavizado & \cite{jm3}, \cite{bishop} \\
4 & Logits, softmax y generación & \cite{d2l}, \cite{gbc}, \cite{holtzman} \\
5 & Entropía cruzada y perplejidad & \cite{jm3}, \cite{ct}, \cite{mackay} \\
\bottomrule
\end{tabular}
\end{center}

Cada tema termina con una lista de lecturas y con ejercicios graduados en tres niveles: $\star$ comprensión, $\star\star$ cálculo y $\star\star\star$ demostración. Los ejercicios marcados con \emph{Fuente} están adaptados de cursos y libros de texto universitarios; el enunciado se ha redactado de nuevo y, en algunos casos, se han cambiado los datos para que puedan resolverse a mano. Las soluciones completas, paso a paso, están en el \cref{ap:soluciones}.

\section*{Sobre las fuentes}
El texto es de redacción propia. Sigue la estructura y la notación del capítulo~3 de Jurafsky y Martin~\cite{jm3}, libro de texto de los cursos de procesamiento del lenguaje natural de Stanford, y de la clase~4 del curso CS224N~\cite{cs224n}. Los resultados de probabilidad siguen a Bertsekas y Tsitsiklis~\cite{bt}, texto del curso 6.041 del MIT~\cite{mit6041}, y los de teoría de la información a Cover y Thomas~\cite{ct}.

\section*{Notación}
\begin{center}
\begin{tabular}{@{}ll@{}}
\toprule
Símbolo & Significado \\
\midrule
$\V$, $K=\abs{\V}$ & vocabulario y su tamaño \\
$\V^{*}$ & conjunto de secuencias finitas de elementos de $\V$ \\
$x_{1:T}=(x_1,\dots,x_T)$ & secuencia de $T$ tokens \\
\fin & símbolo de inicio y fin de secuencia \\
$N(u,v)$ & número de apariciones del bigrama $(u,v)$ en el corpus \\
$z\in\R^{K}$, $T>0$ & vector de logits y temperatura \\
$H(p)$, $H(p,q)$ & entropía y entropía cruzada \\
$\ln$, $\log_2$ & logaritmo natural (nats) y en base 2 (bits) \\
\bottomrule
\end{tabular}
\end{center}

\tableofcontents

\mainmatter
% =====================================================================
\chapter{Vocabulario, tokenización y modelos de lenguaje}
\label{tema:1}
% =====================================================================

Este tema fija los objetos sobre los que trabaja todo el curso: el vocabulario, las secuencias de tokens y la distribución de probabilidad que llamamos modelo de lenguaje. La presentación sigue la introducción del capítulo~3 de~\cite{jm3} y la clase~4 de~\cite{cs224n}.

\section{Vocabulario y secuencias}

\begin{definicion}[Vocabulario]
Un \emph{vocabulario} es un conjunto finito no vacío $\V$ cuyos elementos se llaman \emph{tokens}. Se escribe $K=\abs{\V}$ y se fija una biyección $\V\to\{0,1,\dots,K-1\}$ que asigna a cada token su \emph{índice}.
\end{definicion}

\begin{definicion}[Secuencias]
Para $T\ge 0$, $\V^{T}$ denota el conjunto de secuencias $x_{1:T}=(x_1,\dots,x_T)$ con $x_t\in\V$; $\V^{0}$ contiene solo la secuencia vacía. El conjunto de todas las secuencias finitas es
\[
\V^{*}=\bigcup_{T\ge 0}\V^{T}.
\]
\end{definicion}

\begin{proposicion}\label{prop:contar-secuencias}
Si $K=\abs{\V}\ge 2$, el número de secuencias de longitud a lo sumo $L$ es
\[
\sum_{T=0}^{L} K^{T}=\frac{K^{L+1}-1}{K-1}.
\]
En particular, $\V^{*}$ es infinito numerable.
\end{proposicion}

\begin{proof}
Cada posición de una secuencia de longitud $T$ admite $K$ elecciones independientes, luego $\abs{\V^{T}}=K^{T}$. La suma es una progresión geométrica de razón $K\neq 1$, cuyo valor es el indicado. Como $\V^{*}$ es unión numerable de conjuntos finitos no vacíos, es infinito numerable.
\end{proof}

\section{Tokenización}

El texto se representa como una cadena de caracteres de un alfabeto $\Sigma$ (por ejemplo, Unicode). Un modelo de lenguaje no opera sobre $\Sigma^{*}$, sino sobre secuencias de tokens.

\begin{definicion}[Tokenizador]
Un \emph{tokenizador} es una aplicación $\tau\colon\Sigma^{*}\to\V^{*}$ junto con una aplicación de \emph{decodificación} $\delta\colon\V^{*}\to\Sigma^{*}$ tal que $\delta(\tau(s))=s$ para toda cadena $s$. La condición implica que $\tau$ es inyectiva: dos textos distintos nunca producen la misma secuencia de tokens.
\end{definicion}

\begin{ejemplo}[Tokenización por caracteres]\label{ej:luna}
Sea $\V=\{\fin,\texttt{a},\texttt{l},\texttt{n},\texttt{u}\}$ con índices $0,\dots,4$ en ese orden, donde \fin\ es un token reservado que marca el inicio y el fin de cada palabra. La palabra «luna» se tokeniza como
\[
\tau(\text{luna})=(\fin,\texttt{l},\texttt{u},\texttt{n},\texttt{a},\fin)\ \longmapsto\ (0,2,4,3,1,0).
\]
\end{ejemplo}

La elección de $\V$ determina la granularidad. El \cref{tab:granularidad} resume las tres opciones habituales; los modelos actuales usan subpalabras obtenidas con el algoritmo BPE, que se estudia en el módulo~12.

\begin{table}[h]
\centering\small
\caption{Granularidad de la tokenización.}\label{tab:granularidad}
\begin{tabular}{@{}lccl@{}}
\toprule
Unidad & $K$ típico & Longitud de la secuencia & Limitación principal \\
\midrule
Carácter & $\sim 10^{2}$ & muy larga & secuencias costosas de procesar \\
Palabra & $>5\cdot 10^{5}$ & corta & palabras fuera del vocabulario \\
Subpalabra (BPE) & $3\cdot 10^{4}$–$2{,}6\cdot 10^{5}$ & intermedia & requiere entrenar el tokenizador \\
\bottomrule
\end{tabular}
\end{table}

\section{Modelos de lenguaje}

\begin{definicion}[Modelo de lenguaje]\label{def:ml}
Un \emph{modelo de lenguaje} sobre $\V$ es una distribución de probabilidad $P$ sobre $\V^{*}$, es decir, una función $P\colon\V^{*}\to[0,1]$ tal que
\[
\sum_{x\in\V^{*}}P(x)=1.
\]
\end{definicion}

La suma es una serie, porque $\V^{*}$ es infinito (\cref{prop:contar-secuencias}). Que una construcción concreta cumpla la condición de normalización no es automático: el \cref{tema:2} muestra que un modelo de bigramas sin símbolo de fin la incumple.

\begin{observacion}
En la práctica no se trabaja con $P(x)$ directamente, sino con las probabilidades condicionadas del token siguiente, $P(x_t\mid x_{1:t-1})$. Un gran modelo de lenguaje es una función con parámetros $\theta$ que recibe el contexto $x_{1:t-1}$ y devuelve un vector de $K$ números reales (los \emph{logits}); la función softmax (\cref{tema:4}) los convierte en una distribución sobre $\V$. La regla de la cadena (\cref{tema:2}) conecta ambas descripciones.
\end{observacion}

\begin{lecturas}
\item Jurafsky y Martin~\cite{jm3}, introducción del cap.~3 «N-gram Language Models».
\item Stanford CS224N~\cite{cs224n}, clase~4 «Language Models and RNNs», primera parte.
\item Zhang \emph{et al.}~\cite{d2l}, §9.3 «Language Models», introducción.
\end{lecturas}

\section*{Ejercicios del tema 1}
\addcontentsline{toc}{section}{Ejercicios del tema 1}

\begin{ejercicio}{1}{eje:1-contar}
Sea $\V=\{\fin,\texttt{a},\texttt{l},\texttt{n},\texttt{u}\}$. ¿Cuántas secuencias de longitud a lo sumo~3 hay? ¿Cuántas de longitud exactamente~3?
\end{ejercicio}

\begin{ejercicio}{1}{eje:1-inyectiva}
Demuestre que la condición $\delta(\tau(s))=s$ implica que $\tau$ es inyectiva. ¿Implica también que $\delta$ es inyectiva? Justifique la respuesta.
\end{ejercicio}

\begin{ejercicio}{1}{eje:1-granularidad}
Un texto en español de $10^{4}$ caracteres se tokeniza por caracteres y con un tokenizador de subpalabras que produce, de media, $3{,}5$ caracteres por token. ¿Cuántos tokens resultan en cada caso? Explique qué ventaja aporta la tokenización por subpalabras a un modelo que procesa contextos de longitud fija.
\end{ejercicio}

% =====================================================================
\chapter{Regla de la cadena y modelos de Markov}
\label{tema:2}
% =====================================================================

Este tema descompone la probabilidad de una secuencia en factores de la forma $P(x_t\mid x_{1:t-1})$ y presenta la aproximación de Markov que da lugar a los modelos de $n$-gramas. Los resultados de probabilidad siguen a~\cite[§1.3]{bt}; la parte de $n$-gramas, a~\cite[§3.1]{jm3}.

\section{Probabilidad condicionada}

\begin{definicion}[Probabilidad condicionada]
Sea $(\Omega,\mathcal{F},P)$ un espacio de probabilidad y sean $A,B\in\mathcal{F}$ con $P(A)>0$. La \emph{probabilidad de $B$ condicionada a $A$} es
\[
P(B\mid A)=\frac{P(A\cap B)}{P(A)}.
\]
\end{definicion}

Para $A$ fijo, la aplicación $B\mapsto P(B\mid A)$ es a su vez una medida de probabilidad sobre $(\Omega,\mathcal{F})$~\cite[§1.3]{bt}. Despejando de la definición se obtiene la \emph{regla del producto}, $P(A\cap B)=P(A)\,P(B\mid A)$.

\section{La regla de la cadena}

\begin{teorema}[Regla de la cadena]\label{teo:cadena}
Sean $A_1,\dots,A_T$ sucesos con $P(A_1\cap\dots\cap A_{T-1})>0$. Entonces
\[
P\Bigl(\,\bigcap_{t=1}^{T}A_t\Bigr)=\prod_{t=1}^{T}P\bigl(A_t\mid A_1\cap\dots\cap A_{t-1}\bigr),
\]
donde el factor con $t=1$ se interpreta como $P(A_1)$.
\end{teorema}

\begin{proof}
Por inducción sobre $T$. Para $T=1$ ambos miembros valen $P(A_1)$. Supóngase cierto para $T-1$ y sea $B=A_1\cap\dots\cap A_{T-1}$. Como $P(B)>0$, la regla del producto da
$P(B\cap A_T)=P(B)\,P(A_T\mid B)$. Por la hipótesis de inducción (aplicable porque $P(A_1\cap\dots\cap A_{t-1})\ge P(B)>0$ para todo $t\le T-1$), $P(B)=\prod_{t=1}^{T-1}P(A_t\mid A_1\cap\dots\cap A_{t-1})$. Sustituyendo se obtiene el enunciado.
\end{proof}

\begin{corolario}\label{cor:cadena-secuencias}
Para un modelo de lenguaje y una secuencia $x_{1:T}$ con $P(x_{1:T-1})>0$,
\[
P(x_{1:T})=\prod_{t=1}^{T}P(x_t\mid x_{1:t-1}).
\]
\end{corolario}

\begin{proof}
Basta aplicar el \cref{teo:cadena} a los sucesos $A_t=\{\text{el token en la posición } t \text{ es } x_t\}$.
\end{proof}

La igualdad del \cref{cor:cadena-secuencias} es exacta. Todo modelo de lenguaje autorregresivo, desde un modelo de bigramas hasta un Transformer, consiste en una forma concreta de calcular cada factor $P(x_t\mid x_{1:t-1})$.

\section{Hipótesis de Markov y modelos de \texorpdfstring{$n$}{n}-gramas}

El contexto $x_{1:t-1}$ crece con $t$, y el número de contextos posibles crece exponencialmente (\cref{prop:contar-secuencias}). La solución clásica consiste en truncar el contexto~\cite[§3.1.1]{jm3}.

\begin{definicion}[Modelo de $n$-gramas]\label{def:ngramas}
Sea $n\ge 1$. Un modelo de lenguaje satisface la \emph{hipótesis de Markov de orden $n-1$} si
\[
P(x_t\mid x_{1:t-1})=P(x_t\mid x_{t-n+1:t-1})\qquad\text{para todo } t,
\]
es decir, si la distribución del token siguiente depende solo de los $n-1$ tokens anteriores. El modelo se llama de \emph{unigramas} ($n=1$), de \emph{bigramas} ($n=2$) o de \emph{trigramas} ($n=3$).
\end{definicion}

\begin{definicion}[Matriz estocástica]
Una matriz $P\in\R^{K\times K}$ es \emph{estocástica por filas} si $P_{ij}\ge 0$ para todo $i,j$ y $\sum_{j}P_{ij}=1$ para todo $i$. Un modelo de bigramas queda determinado por la matriz estocástica $P_{ij}=P(x_t=j\mid x_{t-1}=i)$, llamada \emph{matriz de transición}.
\end{definicion}

\begin{proposicion}\label{prop:parametros}
Un modelo de $n$-gramas sobre un vocabulario de tamaño $K$ se especifica con una tabla de $K^{n}$ probabilidades, de las cuales $K^{n-1}(K-1)$ son libres.
\end{proposicion}

\begin{proof}
Hay $K^{n-1}$ contextos posibles y, para cada uno, una distribución sobre $K$ tokens: $K^{n-1}\cdot K=K^{n}$ entradas. Cada distribución está sujeta a una restricción de suma igual a~1, que fija una entrada a partir de las demás; quedan $K-1$ libres por contexto.
\end{proof}

Con $K=5\cdot 10^{4}$ y $n=3$, la tabla tendría $1{,}25\cdot 10^{14}$ entradas. Esta explosión combinatoria es la razón por la que los modelos modernos sustituyen la tabla por una función con parámetros compartidos (módulos~9 y~13).

\begin{ejemplo}[Un modelo de bigramas]\label{ej:bigramas}
Con el vocabulario del \cref{ej:luna} y el corpus «ana», «ala», «luna», el estimador de máxima verosimilitud del \cref{tema:3} da la matriz de transición de la \cref{fig:cadena}. Por el \cref{cor:cadena-secuencias},
\[
P(\text{ana})=P(\texttt{a}\mid\fin)\,P(\texttt{n}\mid\texttt{a})\,P(\texttt{a}\mid\texttt{n})\,P(\fin\mid\texttt{a})=\tfrac{2}{3}\cdot\tfrac{1}{5}\cdot 1\cdot\tfrac{3}{5}=\tfrac{2}{25}=0{,}08.
\]
Del mismo modo, $P(\text{luna})=0{,}10$, $P(\text{ala})=0{,}04$ y $P(\text{lana})=0{,}02$. La palabra «lana» no aparece en el corpus y recibe probabilidad positiva: el modelo combina transiciones observadas.
\end{ejemplo}

\begin{figure}[h]
\centering
\begin{tikzpicture}[->,>={Stealth[length=2mm]},node distance=2.4cm,every state/.style={minimum size=9mm,draw=black,thick},auto,font=\small]
  \node[state] (f) {\fin};
  \node[state] (a) [right=of f] {\texttt{a}};
  \node[state] (n) [above right=1.1cm and 2.0cm of a] {\texttt{n}};
  \node[state] (l) [below right=1.1cm and 2.0cm of a] {\texttt{l}};
  \node[state] (u) [right=3.2cm of a] {\texttt{u}};
  \path (f) edge[bend left=18] node {$2/3$} (a)
        (f) edge[bend right=40] node[below] {$1/3$} (l)
        (a) edge[bend left=18] node {$3/5$} (f)
        (a) edge[bend left=10] node {$1/5$} (n)
        (a) edge[bend left=10] node {$1/5$} (l)
        (n) edge[bend left=10] node {$1$} (a)
        (l) edge[bend left=10] node {$1/2$} (a)
        (l) edge node[right] {$1/2$} (u)
        (u) edge node[right] {$1$} (n);
\end{tikzpicture}
\caption{Diagrama de transición del modelo de bigramas del \cref{ej:bigramas}. Cada arco $i\to j$ lleva $P(j\mid i)$; las probabilidades que salen de cada estado suman~1.}
\label{fig:cadena}
\end{figure}

\section{El símbolo de fin y la normalización}

La \cref{def:ml} exige que las probabilidades de \emph{todas} las secuencias finitas sumen~1. El símbolo \fin\ de fin de secuencia es imprescindible para ello~\cite[§3.1, ejercicio 3.5]{jm3}.

\begin{proposicion}\label{prop:longitud-fija}
Sea un modelo de bigramas sin símbolo de fin, con distribución inicial $\pi$ sobre $\V$ y matriz de transición estocástica $P$. Para cada longitud $L\ge 1$,
\[
\sum_{x_{1:L}\in\V^{L}}\pi(x_1)\prod_{t=2}^{L}P_{x_{t-1}x_t}=1.
\]
En consecuencia, la suma sobre todas las longitudes diverge y el modelo no define una distribución sobre $\V^{*}$.
\end{proposicion}

\begin{proof}
Por inducción sobre $L$. Para $L=1$ la suma es $\sum_{x_1}\pi(x_1)=1$. Para pasar de $L$ a $L+1$, se suma primero sobre $x_{L+1}$: el factor nuevo es $\sum_{x_{L+1}}P_{x_Lx_{L+1}}=1$ porque $P$ es estocástica, y queda la suma de longitud $L$, igual a~1. Sumando sobre $L=1,2,\dots$ se obtiene $1+1+\dots=\infty$.
\end{proof}

\begin{observacion}
Con el símbolo de fin, la generación se detiene la primera vez que aparece \fin, y cada longitud recibe solo una parte de la masa total. Si desde cualquier token se alcanza \fin\ con probabilidad positiva, la masa que queda en secuencias de longitud $L$ decrece geométricamente con $L$ y la suma total es~1; el argumento completo es el de las cadenas de Markov absorbentes~\cite[cap.~7]{bt}.
\end{observacion}

\begin{lecturas}
\item Bertsekas y Tsitsiklis~\cite{bt}, §1.3 «Conditional Probability», regla de la multiplicación.
\item MIT 6.041SC~\cite{mit6041}, unidad~I, clase~2 (probabilidad condicionada).
\item Jurafsky y Martin~\cite{jm3}, §3.1 «N-Grams» y §3.1.1 «The Markov assumption».
\item Zhang \emph{et al.}~\cite{d2l}, §9.3.1.1 «Markov Models and $n$-grams».
\end{lecturas}

\section*{Ejercicios del tema 2}
\addcontentsline{toc}{section}{Ejercicios del tema 2}

\begin{ejercicio}{1}{eje:2-casillas}
Con un vocabulario de 27 tokens (26 letras y \fin), ¿cuántas entradas tiene la tabla de un modelo de bigramas? ¿Y la de uno de trigramas? ¿Cuántas de ellas son parámetros libres en cada caso?
\end{ejercicio}

\begin{ejercicio}{1}{eje:2-vf}
Justifique si son verdaderas o falsas: (a) la regla de la cadena es una aproximación; (b) cada fila de una matriz de transición suma~1.
\end{ejercicio}

\begin{ejercicio}{2}{eje:2-sol-matriz}
Corpus: «sol», «sal», «los». Con $\V=\{\fin,\texttt{a},\texttt{l},\texttt{o},\texttt{s}\}$, construya la matriz de conteos $N$ y la matriz de transición $P$ del modelo de bigramas (estimadas por frecuencia relativa).
\end{ejercicio}

\begin{ejercicio}{2}{eje:2-sol-prob}
Con el modelo del \cref{eje:2-sol-matriz}, calcule $P(\text{sol})$, $P(\text{sal})$, $P(\text{los})$ y $P(\text{sos})$. Explique por qué $P(\text{sal})=2P(\text{sol})$.
\end{ejercicio}

\begin{ejercicio}{1}{eje:2-longitud}
En el \cref{ej:bigramas}, la palabra más probable es «a», con $P(\text{a})=0{,}40$. Explique por qué y qué problema plantea para comparar palabras de distinta longitud. Proponga una medida que lo corrija y aplíquela a «a», «ana» y «luna».
\end{ejercicio}

\begin{ejercicio}{3}{eje:2-sin-fin}
Se entrena un modelo de bigramas sin suavizar y sin símbolo de fin, con símbolo de inicio $\langle s\rangle$, sobre las frases «$\langle s\rangle$ a b», «$\langle s\rangle$ b b», «$\langle s\rangle$ b a» y «$\langle s\rangle$ a b». (a) Escriba todas las probabilidades de transición. (b) Calcule la probabilidad de las cuatro frases de dos palabras sobre $\{\texttt{a},\texttt{b}\}$ y compruebe que suman~1. (c) Relacione el resultado con la \cref{prop:longitud-fija}.
\fuente{adaptado de Jurafsky y Martin~\cite{jm3}, cap.~3, ejercicio~3.5 (corpus modificado).}
\end{ejercicio}

% =====================================================================
\chapter{Estimación por máxima verosimilitud y suavizado}
\label{tema:3}
% =====================================================================

Los parámetros de un modelo de $n$-gramas se estiman a partir de un corpus. Este tema demuestra que el estimador de máxima verosimilitud es la frecuencia relativa, analiza su principal defecto (las probabilidades nulas) y presenta dos remedios: el suavizado aditivo y la interpolación. Se sigue~\cite[§3.1.2 y §3.6]{jm3} y, para la interpretación bayesiana, \cite[§2.2]{bishop}.

\section{Verosimilitud de un corpus}

Sea un corpus formado por secuencias, cada una delimitada por \fin\ al principio y al final, y sea $N(u,v)$ el número de veces que el token $v$ sigue inmediatamente al token $u$. Se escribe $N(u)=\sum_{v}N(u,v)$ para el número de veces que $u$ aparece como token anterior.

\begin{definicion}[Verosimilitud]
Dada una matriz de transición $P$, la \emph{verosimilitud} del corpus es el producto de las probabilidades de todos sus bigramas,
\[
\mathcal{L}(P)=\prod_{u,v\in\V}P_{uv}^{\,N(u,v)},
\qquad
\ell(P)=\ln\mathcal{L}(P)=\sum_{u,v\in\V}N(u,v)\ln P_{uv},
\]
y $\ell$ se llama \emph{log-verosimilitud}. El \emph{estimador de máxima verosimilitud} (EMV) es la matriz estocástica que maximiza $\ell$.
\end{definicion}

El logaritmo convierte productos en sumas y evita el subdesbordamiento numérico: con $10^{3}$ factores iguales a $0{,}9$ el producto vale aproximadamente $10^{-46}$.

\section{El estimador de máxima verosimilitud}

\begin{teorema}[EMV del modelo de bigramas]\label{teo:emv}
Supóngase $N(u)>0$ para todo $u$. El estimador de máxima verosimilitud de la matriz de transición es la frecuencia relativa
\[
\widehat{P}_{uv}=\frac{N(u,v)}{N(u)}.
\]
\end{teorema}

\begin{proof}
\emph{Paso 1 (separación por filas).} La función $\ell(P)=\sum_{u}\bigl(\sum_{v}N(u,v)\ln P_{uv}\bigr)$ es suma de términos que dependen cada uno de una sola fila, y las restricciones $\sum_{v}P_{uv}=1$ también afectan a una sola fila. Por tanto, basta maximizar cada fila por separado.

\emph{Paso 2 (lagrangiana).} Fijada una fila $u$, se escribe $N_v=N(u,v)$, $p_v=P_{uv}$ y $n=N(u)$. Se maximiza $f(p)=\sum_{v}N_v\ln p_v$ con la restricción $\sum_v p_v=1$. Si algún $N_v>0$ tuviera $p_v=0$, entonces $f=-\infty$; luego en el máximo $p_v>0$ siempre que $N_v>0$. La lagrangiana es
\[
\Lambda(p,\lambda)=\sum_{v}N_v\ln p_v-\lambda\Bigl(\sum_{v}p_v-1\Bigr).
\]

\emph{Paso 3 (condiciones de primer orden).} $\partial\Lambda/\partial p_v=N_v/p_v-\lambda=0$, de donde $p_v=N_v/\lambda$ para los $v$ con $N_v>0$. Los $v$ con $N_v=0$ no contribuyen a $f$, y cualquier masa asignada a ellos se resta a los demás; por tanto, en el óptimo, $p_v=0$ para ellos.

\emph{Paso 4 (multiplicador).} Imponiendo la restricción, $\sum_{v}N_v/\lambda=1$, luego $\lambda=\sum_v N_v=n$ y $p_v=N_v/n$.

\emph{Paso 5 (máximo global).} $f$ es una suma de funciones cóncavas $N_v\ln p_v$ con $N_v\ge 0$, y el conjunto de distribuciones es convexo; un punto crítico de una función cóncava sobre un convexo es un máximo global.
\end{proof}

\begin{ejemplo}\label{ej:conteos}
Para el corpus «ana», «ala», «luna» los bigramas son 13 y la matriz de conteos es la del \cref{tab:conteos}. Dividiendo cada fila por su total se obtienen las probabilidades de la \cref{fig:cadena}.
\end{ejemplo}

\begin{table}[h]
\centering\small
\caption{Conteos $N(u,v)$ del \cref{ej:conteos} (fila: token anterior; columna: token siguiente).}\label{tab:conteos}
\begin{tabular}{@{}c|ccccc|c@{}}
\toprule
$u\backslash v$ & \fin & \texttt{a} & \texttt{l} & \texttt{n} & \texttt{u} & $N(u)$ \\
\midrule
\fin & 0 & 2 & 1 & 0 & 0 & 3 \\
\texttt{a} & 3 & 0 & 1 & 1 & 0 & 5 \\
\texttt{l} & 0 & 1 & 0 & 0 & 1 & 2 \\
\texttt{n} & 0 & 2 & 0 & 0 & 0 & 2 \\
\texttt{u} & 0 & 0 & 0 & 1 & 0 & 1 \\
\bottomrule
\end{tabular}
\end{table}

\section{Ceros y suavizado aditivo}

El EMV asigna probabilidad nula a todo bigrama que no aparece en el corpus. Si un texto de evaluación contiene uno de ellos, su probabilidad es~0 y la pérdida del \cref{tema:5} es infinita. Este fenómeno es una forma de sobreajuste~\cite[§3.5]{jm3}.

\begin{definicion}[Suavizado aditivo]\label{def:laplace}
Sea $\alpha>0$. El estimador con \emph{suavizado aditivo} es
\[
\widehat{P}_{\alpha}(v\mid u)=\frac{N(u,v)+\alpha}{N(u)+\alpha K}.
\]
Con $\alpha=1$ recibe el nombre de \emph{suavizado de Laplace}; con $0<\alpha<1$, de \emph{suavizado add-$k$}~\cite[§3.6.1–3.6.2]{jm3}.
\end{definicion}

\begin{proposicion}\label{prop:laplace}
Para cada $u$, $\widehat{P}_{\alpha}(\cdot\mid u)$ es una distribución de probabilidad con todas sus entradas estrictamente positivas. Además, coincide con la media a posteriori de la fila $u$ cuando se adopta una distribución a priori de Dirichlet de parámetros $(\alpha,\dots,\alpha)$.
\end{proposicion}

\begin{proof}
Positividad: el numerador es al menos $\alpha>0$. Suma: $\sum_{v}(N(u,v)+\alpha)=N(u)+\alpha K$, que es el denominador. Para la segunda afirmación, la distribución de Dirichlet es conjugada de la multinomial: si la fila $u$ tiene a priori $\mathrm{Dir}(\alpha,\dots,\alpha)$ y se observan los conteos $N(u,\cdot)$, la distribución a posteriori es $\mathrm{Dir}(\alpha+N(u,v_1),\dots,\alpha+N(u,v_K))$, cuya media en la coordenada $v$ es $(N(u,v)+\alpha)/(N(u)+\alpha K)$~\cite[§2.2.1]{bishop}.
\end{proof}

\begin{ejemplo}
En el \cref{ej:conteos} con $\alpha=1$ y $K=5$: $\widehat{P}_1(\texttt{n}\mid\fin)=(0+1)/(3+5)=1/8$, mientras que el EMV daba~0; y $\widehat{P}_1(\texttt{a}\mid\texttt{n})=(2+1)/(2+5)=3/7\approx 0{,}4286$, frente al valor~1 del EMV. El suavizado transfiere masa de los sucesos observados a los no observados.
\end{ejemplo}

\section{Interpolación}

Otra estrategia consiste en combinar estimadores de distinto orden~\cite[§3.6.3]{jm3}.

\begin{definicion}[Interpolación lineal]
Dados un estimador de bigramas $\widehat{P}_2(v\mid u)$ y uno de unigramas $\widehat{P}_1(v)$, y pesos $\lambda_1,\lambda_2\ge 0$ con $\lambda_1+\lambda_2=1$, el estimador \emph{interpolado} es
\[
\widehat{P}_{\mathrm{int}}(v\mid u)=\lambda_2\,\widehat{P}_2(v\mid u)+\lambda_1\,\widehat{P}_1(v).
\]
\end{definicion}

Es una distribución, por ser combinación convexa de distribuciones. Los pesos se eligen maximizando la verosimilitud de un conjunto de validación separado del de entrenamiento.

\begin{lecturas}
\item Jurafsky y Martin~\cite{jm3}, §3.1.2 «How to estimate probabilities», §3.5 «Generalizing vs.\ overfitting the training set» y §3.6 «Smoothing, Interpolation, and Backoff».
\item Bishop~\cite{bishop}, §2.2 «Multinomial variables» y §2.2.1 «The Dirichlet distribution».
\item Chen y Goodman~\cite{chen-goodman}, estudio comparativo de técnicas de suavizado.
\end{lecturas}

\section*{Ejercicios del tema 3}
\addcontentsline{toc}{section}{Ejercicios del tema 3}

\begin{ejercicio}{2}{eje:3-laplace}
Aplique el suavizado de Laplace ($\alpha=1$) al modelo del \cref{eje:2-sol-matriz}. Calcule $P(\fin\mid\fin)$ y $P(\texttt{l}\mid\texttt{a})$ y compárelas con las del EMV.
\end{ejercicio}

\begin{ejercicio}{2}{eje:3-u1}
Corpus de cuatro frases con marcadores $\langle s\rangle$ y $\langle/s\rangle$, que cuentan como tokens del vocabulario: «$\langle s\rangle$ yo soy Ana $\langle/s\rangle$», «$\langle s\rangle$ Ana yo soy $\langle/s\rangle$», «$\langle s\rangle$ yo soy Ana $\langle/s\rangle$» y «$\langle s\rangle$ yo no quiero sopa ni Ana $\langle/s\rangle$». Calcule $P(\text{Ana}\mid\text{soy})$ con un modelo de bigramas suavizado con Laplace.
\fuente{adaptado de Jurafsky y Martin~\cite{jm3}, cap.~3, ejercicio~3.4.}
\end{ejercicio}

\begin{ejercicio}{2}{eje:3-interp}
Con el corpus del \cref{eje:3-u1}, calcule $P(\text{Ana}\mid\text{soy})$ interpolando el bigrama y el unigrama de máxima verosimilitud con $\lambda_1=\lambda_2=1/2$. Para el unigrama, cuente todos los tokens, incluidos $\langle s\rangle$ y $\langle/s\rangle$.
\fuente{adaptado de Jurafsky y Martin~\cite{jm3}, cap.~3, ejercicio~3.7.}
\end{ejercicio}

\begin{ejercicio}{3}{eje:3-limite}
Demuestre que, para $u$ fijo con $N(u)>0$, $\lim_{\alpha\to 0^{+}}\widehat{P}_{\alpha}(v\mid u)$ es el EMV y $\lim_{\alpha\to\infty}\widehat{P}_{\alpha}(v\mid u)=1/K$. Interprete $\alpha$ como un número de «observaciones ficticias».
\end{ejercicio}

% =====================================================================
\chapter{Logits, softmax y generación}
\label{tema:4}
% =====================================================================

Una red neuronal no produce probabilidades directamente, sino un vector de números reales sin restricciones. Este tema estudia la función softmax, que transforma ese vector en una distribución, su versión con temperatura y los procedimientos de muestreo con los que un modelo genera texto. Se sigue~\cite[§4.1]{d2l}, \cite[§4.1]{gbc} y~\cite[§3.4]{jm3}.

\section{La función softmax}

\begin{definicion}[Softmax con temperatura]\label{def:softmax}
Sea $z\in\R^{K}$ un vector de \emph{logits} y $T>0$ una \emph{temperatura}. Se define
\[
\softmax_T(z)_i=\frac{e^{z_i/T}}{\sum_{j=1}^{K}e^{z_j/T}},\qquad i=1,\dots,K.
\]
Con $T=1$ se escribe simplemente $\softmax(z)$.
\end{definicion}

\begin{proposicion}[Propiedades]\label{prop:softmax}
Sean $z\in\R^{K}$ y $T>0$.
\begin{enumerate}
  \item $\softmax_T(z)$ pertenece al símplex: sus componentes son positivas y suman~1.
  \item Invariancia por traslación: $\softmax_T(z+c\mathbf{1})=\softmax_T(z)$ para todo $c\in\R$.
  \item Monotonía: si $z_i>z_j$, entonces $\softmax_T(z)_i>\softmax_T(z)_j$.
  \item Si el máximo $z_k$ es único, $\lim_{T\to 0^{+}}\softmax_T(z)=e_k$, el vector con un~1 en la posición $k$.
  \item $\lim_{T\to\infty}\softmax_T(z)=(1/K,\dots,1/K)$.
\end{enumerate}
\end{proposicion}

\begin{proof}
(i) Cada término $e^{z_i/T}$ es positivo y el denominador es la suma de todos ellos. (ii) $e^{(z_i+c)/T}=e^{c/T}e^{z_i/T}$ y el factor común $e^{c/T}>0$ se cancela entre numerador y denominador. (iii) La función $s\mapsto e^{s/T}$ es estrictamente creciente y el denominador es común. (iv) Por (ii) con $c=-z_k$, $\softmax_T(z)_i=e^{(z_i-z_k)/T}/\sum_j e^{(z_j-z_k)/T}$. El término $j=k$ vale~1; para $j\ne k$, $z_j-z_k<0$ y $e^{(z_j-z_k)/T}\to 0$. El denominador tiende a~1, luego la componente $k$ tiende a~1 y las demás a~0. (v) Cada $z_j/T\to 0$ y, por continuidad de la exponencial, cada término tiende a~1; el cociente tiende a $1/K$.
\end{proof}

El \cref{tab:temperatura} ilustra la proposición para $z=(2,1,0)$. La temperatura no altera el orden de los tokens (propiedad~(iii)); solo controla cuánto se concentra la distribución.

\begin{table}[h]
\centering\small
\caption{$\softmax_T(z)$ para $z=(2,1,0)$.}\label{tab:temperatura}
\begin{tabular}{@{}lccc@{}}
\toprule
Temperatura & $p_1$ & $p_2$ & $p_3$ \\
\midrule
$T\to 0^{+}$ & $1$ & $0$ & $0$ \\
$T=0{,}5$ & $0{,}867$ & $0{,}117$ & $0{,}016$ \\
$T=1$ & $0{,}665$ & $0{,}245$ & $0{,}090$ \\
$T=2$ & $0{,}506$ & $0{,}307$ & $0{,}186$ \\
$T\to\infty$ & $1/3$ & $1/3$ & $1/3$ \\
\bottomrule
\end{tabular}
\end{table}

\begin{observacion}[Relación con los modelos de $n$-gramas]
Si en una fila de un modelo de bigramas se toman como logits $z_v=\ln\widehat{P}(v\mid u)$ para los $v$ con probabilidad positiva, entonces $\softmax(z)=\widehat{P}(\cdot\mid u)$. Elevar las probabilidades a $1/T$ y renormalizar equivale a aplicar la softmax con temperatura $T$.
\end{observacion}

\begin{observacion}[Relación con la física]
La expresión de la \cref{def:softmax} es la distribución de Boltzmann de la mecánica estadística, $p_i\propto e^{-E_i/(k_BT)}$, con energías $E_i=-z_i$ y la temperatura medida en unidades de $k_B$. Los límites (iv) y (v) corresponden al estado fundamental a temperatura nula y a la equiprobabilidad a temperatura infinita.
\end{observacion}

\section{Estabilidad numérica}

Para valores grandes de $z_i$, el cálculo directo de $e^{z_i}$ desborda la aritmética de coma flotante: en doble precisión el mayor número representable es aproximadamente $1{,}8\cdot 10^{308}$, y $e^{710}$ ya lo supera~\cite[§4.1]{gbc}.

\begin{definicion}[Log-suma-exp]
$\LSE(z)=\ln\sum_{i=1}^{K}e^{z_i}$. Se cumple $\ln\softmax(z)_i=z_i-\LSE(z)$.
\end{definicion}

\begin{proposicion}\label{prop:lse}
Sea $m=\max_i z_i$. Entonces $\LSE(z)=m+\LSE(z-m\mathbf{1})$ y
\[
m\le\LSE(z)\le m+\ln K.
\]
\end{proposicion}

\begin{proof}
$\sum_i e^{z_i}=e^{m}\sum_i e^{z_i-m}$; tomando logaritmos se obtiene la igualdad. En la suma $\sum_i e^{z_i-m}$ hay un término igual a~1 y los demás están en $(0,1]$, luego la suma está entre $1$ y $K$, y su logaritmo entre $0$ y $\ln K$.
\end{proof}

Tras restar el máximo, todos los exponentes son $\le 0$ y no se produce desbordamiento. Las bibliotecas numéricas calculan así la softmax y la entropía cruzada.

\section{Muestreo}

Generar texto consiste en aplicar repetidamente el \cref{cor:cadena-secuencias}: se obtiene la distribución $p=P(\cdot\mid x_{1:t-1})$, se elige un token $x_t$ y se añade al contexto, hasta que aparece \fin~\cite[§3.4]{jm3}.

\begin{proposicion}[Muestreo por la inversa de la función de distribución]\label{prop:muestreo}
Sea $p=(p_1,\dots,p_K)$ una distribución, $F_k=\sum_{i\le k}p_i$ con $F_0=0$, y sea $U$ una variable aleatoria uniforme en $[0,1)$. Defínase $X=\min\{k: U<F_k\}$. Entonces $P(X=k)=p_k$ para todo $k$.
\end{proposicion}

\begin{proof}
$X=k$ si y solo si $F_{k-1}\le U<F_k$. Como $U$ es uniforme, la probabilidad de este suceso es la longitud del intervalo, $F_k-F_{k-1}=p_k$.
\end{proof}

\begin{definicion}[Estrategias de decodificación]
Dada $p=P(\cdot\mid x_{1:t-1})$:
\begin{enumerate}
  \item \emph{decodificación voraz}: $x_t=\arg\max_i p_i$ (equivale a $T\to 0^{+}$);
  \item \emph{muestreo top-$k$}: se conservan los $k$ tokens más probables, se renormaliza y se muestrea~\cite{fan};
  \item \emph{muestreo de núcleo} o \emph{top-$p$}: se conserva el menor conjunto de tokens, ordenados por probabilidad decreciente, cuya probabilidad acumulada alcanza un umbral $p_0\in(0,1]$, se renormaliza y se muestrea~\cite{holtzman}.
\end{enumerate}
\end{definicion}

Estas estrategias, combinadas con la temperatura, son los controles que exponen las interfaces de los modelos actuales. Se estudian con más detalle en el módulo~19.

\begin{lecturas}
\item Zhang \emph{et al.}~\cite{d2l}, §4.1 «Softmax Regression».
\item Goodfellow, Bengio y Courville~\cite{gbc}, §4.1 «Overflow and Underflow».
\item Jurafsky y Martin~\cite{jm3}, §3.4 «Sampling sentences from a language model».
\item Holtzman \emph{et al.}~\cite{holtzman}, definición y análisis del muestreo de núcleo.
\end{lecturas}

\section*{Ejercicios del tema 4}
\addcontentsline{toc}{section}{Ejercicios del tema 4}

\begin{ejercicio}{1}{eje:4-orden}
¿Puede un cambio de temperatura modificar cuál es el token más probable? Justifique la respuesta con la \cref{prop:softmax}.
\end{ejercicio}

\begin{ejercicio}{2}{eje:4-softmax}
Calcule $\softmax_T(z)$ para $z=(0,\ln 2,\ln 5)$ con $T=1$ y con $T=2$. ¿Cuál es el límite cuando $T\to 0^{+}$?
\end{ejercicio}

\begin{ejercicio}{2}{eje:4-muestreo}
Una fila del modelo da $p=(0{,}1;\,0{,}5;\,0{,}4)$ para los tokens $x$, $y$, $z$, en ese orden. Aplicando la \cref{prop:muestreo}, ¿qué token se obtiene para $u=0{,}05$, $0{,}35$, $0{,}6$ y $0{,}95$?
\end{ejercicio}

\begin{ejercicio}{2}{eje:4-potencia}
Dada una distribución $P$, se define $Q(i)\propto P(i)^{\beta}$ con $\beta>0$. (a) ¿Qué valor de $\beta$ equivale a duplicar la temperatura y cuál a reducirla a la mitad? (b) Calcule $Q$ para $P=(0{,}5;\,0{,}25;\,0{,}25)$ con $\beta=2$ y con $\beta=1/2$.
\fuente{adaptado de Zhang \emph{et al.}~\cite{d2l}, §4.1, ejercicio~7.}
\end{ejercicio}

\begin{ejercicio}{3}{eje:4-limites}
Demuestre con detalle las propiedades (ii), (iv) y (v) de la \cref{prop:softmax}. Muestre con un ejemplo que (iv) falla si el máximo no es único.
\end{ejercicio}

\begin{ejercicio}{3}{eje:4-lse}
Sea $g(x)=\LSE(x)$. (a) Demuestre que $\partial g/\partial x_i=\softmax(x)_i$. (b) Demuestre que $g(x+b\mathbf{1})=g(x)+b$. (c) Calcule $g(1000;\,1001;\,1002)$ a mano y explique qué ocurre si se evalúa directamente en doble precisión.
\fuente{adaptado de Zhang \emph{et al.}~\cite{d2l}, §4.1, ejercicio~6, apartados (b) y (d).}
\end{ejercicio}

% =====================================================================
\chapter{Evaluación: entropía cruzada y perplejidad}
\label{tema:5}
% =====================================================================

Un buen modelo asigna probabilidad alta al texto real. Este tema define la pérdida que se minimiza al entrenar todos los modelos de lenguaje, la entropía cruzada, y la medida de evaluación derivada de ella, la perplejidad. Se sigue~\cite[§3.2, §3.3 y §3.7]{jm3}; los resultados de teoría de la información proceden de~\cite[cap.~2]{ct}.

\section{Pérdida y perplejidad}

Sea un texto de evaluación con $N$ predicciones y sea $p_t=P(x_t\mid x_{1:t-1})\in(0,1]$ la probabilidad que el modelo asigna al token que realmente aparece en la posición $t$.

\begin{definicion}[Pérdida y perplejidad]\label{def:ppl}
La \emph{pérdida} (log-verosimilitud negativa media) y la \emph{perplejidad} del modelo sobre el texto son
\[
\mathcal{J}=-\frac{1}{N}\sum_{t=1}^{N}\ln p_t,\qquad \PPL=e^{\mathcal{J}}.
\]
Con logaritmo natural, $\mathcal{J}$ se mide en \emph{nats}; con $\log_2$, en \emph{bits}, y $\mathcal{J}_{\text{bits}}=\mathcal{J}_{\text{nats}}/\ln 2$.
\end{definicion}

\begin{proposicion}\label{prop:ppl}
\begin{enumerate}
  \item $\PPL=\bigl(\prod_{t=1}^{N}1/p_t\bigr)^{1/N}$, la media geométrica de los inversos de las probabilidades.
  \item $\PPL\ge 1$, con igualdad si y solo si $p_t=1$ para todo $t$.
  \item Un modelo uniforme sobre $K$ tokens tiene $\PPL=K$ sobre cualquier texto.
\end{enumerate}
\end{proposicion}

\begin{proof}
(i) $-\ln p_t=\ln(1/p_t)$, y la suma de logaritmos es el logaritmo del producto: $\mathcal{J}=\frac1N\ln\prod_t(1/p_t)=\ln\bigl(\prod_t 1/p_t\bigr)^{1/N}$; basta aplicar la exponencial. (ii) Cada $1/p_t\ge 1$; el producto y su raíz $N$-ésima son $\ge 1$, con igualdad si y solo si todos los factores valen~1. (iii) Si $p_t=1/K$ para todo $t$, el producto es $K^{N}$ y su raíz $N$-ésima es $K$.
\end{proof}

La perplejidad admite una interpretación como \emph{factor de ramificación medio}: un modelo con perplejidad $k$ está, en promedio, tan indeciso como si eligiera uniformemente entre $k$ tokens~\cite[§3.3.1]{jm3}.

\begin{ejemplo}\label{ej:ppl}
El modelo del \cref{ej:bigramas}, evaluado sobre su propio corpus (13 bigramas), tiene $\sum_t\ln p_t=-8{,}047$. Por tanto, $\mathcal{J}=8{,}047/13=0{,}619$ nats $=0{,}893$ bits y $\PPL=e^{0{,}619}=1{,}86$. El modelo uniforme sobre los 5 tokens tiene $\mathcal{J}=\ln 5=1{,}609$ nats y $\PPL=5$.
\end{ejemplo}

\section{Entropía y entropía cruzada}

\begin{definicion}[Entropía y entropía cruzada]
Sean $p$ y $q$ distribuciones sobre un conjunto finito $\mathcal{X}$, con $q(x)>0$ siempre que $p(x)>0$. La \emph{entropía} de $p$ y la \emph{entropía cruzada} de $q$ respecto de $p$ son
\[
H(p)=-\sum_{x\in\mathcal{X}}p(x)\ln p(x),\qquad
H(p,q)=-\sum_{x\in\mathcal{X}}p(x)\ln q(x),
\]
con el convenio $0\ln 0=0$.
\end{definicion}

Si los tokens del texto se distribuyen según $p$ y el modelo es $q$, la pérdida $\mathcal{J}$ es un estimador de $H(p,q)$: es la media empírica de $-\ln q(X)$ con $X\sim p$.

\begin{teorema}[Desigualdad de Gibbs]\label{teo:gibbs}
$H(p,q)\ge H(p)$, con igualdad si y solo si $q=p$.
\end{teorema}

\begin{proof}
Sea $S=\{x:p(x)>0\}$. Entonces
\[
H(p)-H(p,q)=\sum_{x\in S}p(x)\ln\frac{q(x)}{p(x)}\le\ln\sum_{x\in S}p(x)\frac{q(x)}{p(x)}=\ln\sum_{x\in S}q(x)\le\ln 1=0,
\]
donde la primera desigualdad es la de Jensen para la función estrictamente cóncava $\ln$. La igualdad en Jensen exige que $q(x)/p(x)$ sea constante en $S$, y la igualdad final exige $\sum_{x\in S}q(x)=1$; ambas juntas dan $q=p$~\cite[teorema~2.6.3]{ct}.
\end{proof}

La diferencia $H(p,q)-H(p)=\sum_x p(x)\ln\bigl(p(x)/q(x)\bigr)$ es la \emph{divergencia de Kullback–Leibler} $D(p\,\|\,q)$, que se estudia en el módulo~6.

\begin{corolario}\label{cor:aprender}
Entre todos los modelos $q$, la entropía cruzada esperada $H(p,q)$ se minimiza exactamente en $q=p$, y su valor mínimo es $H(p)$.
\end{corolario}

El \cref{cor:aprender} justifica el criterio de entrenamiento de los modelos de lenguaje: minimizar la entropía cruzada media sobre un corpus empuja al modelo a reproducir la distribución real del token siguiente. Ningún modelo puede bajar de $H(p)$, la incertidumbre propia del lenguaje. La relación entre entropía y predicción se remonta a Shannon~\cite{shannon}.

\section{Entrenamiento y evaluación}

\begin{observacion}
La perplejidad debe medirse sobre un conjunto de evaluación disjunto del de entrenamiento. Evaluar sobre los datos de entrenamiento da una estimación optimista: el EMV del \cref{teo:emv} maximiza precisamente la verosimilitud de esos datos~\cite[§3.2]{jm3}. Además, la perplejidad solo es comparable entre modelos que usan el mismo vocabulario, porque depende de la tokenización.
\end{observacion}

\begin{lecturas}
\item Jurafsky y Martin~\cite{jm3}, §3.2 «Evaluating Language Models: Training and Test Sets», §3.3 «Evaluating Language Models: Perplexity» y §3.7 «Perplexity's Relation to Entropy».
\item Cover y Thomas~\cite{ct}, §2.1 «Entropy», §2.3 «Relative Entropy and Mutual Information» y §2.6 «Jensen's Inequality and Its Consequences».
\item MacKay~\cite{mackay}, cap.~2 (entropía y desigualdad de Gibbs).
\item Zhang \emph{et al.}~\cite{d2l}, §9.3.2 «Perplexity».
\end{lecturas}

\section*{Ejercicios del tema 5}
\addcontentsline{toc}{section}{Ejercicios del tema 5}

\begin{ejercicio}{1}{eje:5-ppl1}
¿Qué significa que un modelo tenga perplejidad~1 sobre un texto? Justifíquelo con la \cref{prop:ppl}.
\end{ejercicio}

\begin{ejercicio}{2}{eje:5-perdida}
Un modelo asigna las probabilidades $1/2$, $1/4$, $1/4$ y $1/8$ a los cuatro tokens de un texto de evaluación. Calcule la pérdida en bits y en nats, y la perplejidad.
\end{ejercicio}

\begin{ejercicio}{2}{eje:5-digitos}
Un conjunto de entrenamiento contiene 100 dígitos: 82 ceros y dos de cada dígito del 1 al 9. Calcule la perplejidad de un modelo de unigramas de máxima verosimilitud sobre el texto de prueba «0 0 0 0 0 3 0 0 0 0» y compárela con la de un modelo uniforme sobre los diez dígitos.
\fuente{adaptado de Jurafsky y Martin~\cite{jm3}, cap.~3, ejercicio~3.12 (conteos modificados).}
\end{ejercicio}

\begin{ejercicio}{3}{eje:5-ppl}
Demuestre la \cref{prop:ppl} con todo detalle, justificando cada propiedad del logaritmo que utilice.
\end{ejercicio}

\begin{ejercicio}{3}{eje:5-lagrange}
Demuestre el \cref{cor:aprender} sin usar la desigualdad de Jensen, mediante multiplicadores de Lagrange: minimice $H(p,q)$ en $q$ con la restricción $\sum_x q(x)=1$, suponiendo $p(x)>0$ para todo $x$.
\end{ejercicio}


\appendix
% =====================================================================
\chapter{Soluciones de los ejercicios}
\label{ap:soluciones}
% =====================================================================

Cada solución se presenta en cinco partes: datos, planteamiento, desarrollo, resultado y comprobación. En los pasos intermedios se redondea a cuatro cifras significativas salvo que se indique otra cosa; los resultados exactos se dan como fracción cuando existe.

% ---------------------------------------------------------------------
\section{Tema 1}

\begin{solucion}{eje:1-contar}
\datos $K=\abs{\V}=5$; longitudes $T=0,1,2,3$.
\planteamiento por la \cref{prop:contar-secuencias}, hay $K^{T}$ secuencias de longitud $T$.
\pasos
\begin{enumerate}[label=\arabic*.]
  \item $K^{0}+K^{1}+K^{2}+K^{3}=1+5+25+125=156$.
  \item Con la fórmula: $(K^{4}-1)/(K-1)=(625-1)/4=624/4=156$.
  \item Longitud exactamente~3: $K^{3}=125$.
\end{enumerate}
\resultado \textbf{156} secuencias de longitud a lo sumo~3; \textbf{125} de longitud exactamente~3.
\comprobacion la secuencia vacía cuenta como una; $156-125=31=1+5+25$ son las de longitud menor que~3.
\end{solucion}

\begin{solucion}{eje:1-inyectiva}
\datos $\delta(\tau(s))=s$ para toda cadena $s\in\Sigma^{*}$.
\planteamiento para probar la inyectividad de $\tau$, se parte de $\tau(s)=\tau(s')$ y se aplica $\delta$.
\pasos
\begin{enumerate}[label=\arabic*.]
  \item Si $\tau(s)=\tau(s')$, aplicando $\delta$ a ambos lados: $s=\delta(\tau(s))=\delta(\tau(s'))=s'$. Luego $\tau$ es inyectiva.
  \item $\delta$ no tiene por qué ser inyectiva. Contraejemplo: si $\V$ contiene los tokens \texttt{a}, \texttt{b} y \texttt{ab}, las secuencias $(\texttt{a},\texttt{b})$ y $(\texttt{ab})$ son distintas y ambas se decodifican como la cadena «ab».
\end{enumerate}
\resultado \textbf{$\tau$ es inyectiva; $\delta$ en general no lo es}, aunque sí es sobreyectiva.
\comprobacion la sobreyectividad de $\delta$ se sigue de la misma condición: toda cadena $s$ es la imagen de $\tau(s)$.
\end{solucion}

\begin{solucion}{eje:1-granularidad}
\datos $10^{4}$ caracteres; $3{,}5$ caracteres por token en el tokenizador de subpalabras.
\planteamiento número de tokens $=$ número de caracteres $/$ caracteres por token.
\pasos
\begin{enumerate}[label=\arabic*.]
  \item Por caracteres: un token por carácter, $10^{4}$ tokens.
  \item Por subpalabras: $10^{4}/3{,}5=2857{,}1\approx 2857$ tokens.
  \item Con un contexto fijo de $C$ tokens, el tokenizador de subpalabras cubre $3{,}5\,C$ caracteres de texto, frente a $C$ caracteres por caracteres.
\end{enumerate}
\resultado \textbf{$10^{4}$ tokens frente a unos $2857$}. La tokenización por subpalabras permite que un contexto de longitud fija abarque unas $3{,}5$ veces más texto.
\comprobacion $2857\cdot 3{,}5=9999{,}5\approx 10^{4}$.
\end{solucion}

% ---------------------------------------------------------------------
\section{Tema 2}

\begin{solucion}{eje:2-casillas}
\datos $K=27$.
\planteamiento por la \cref{prop:parametros}, la tabla de $n$-gramas tiene $K^{n}$ entradas y $K^{n-1}(K-1)$ parámetros libres.
\pasos
\begin{enumerate}[label=\arabic*.]
  \item Bigramas: $K^{2}=27\cdot 27=729$ entradas; libres $27\cdot 26=702$.
  \item Trigramas: $K^{3}=729\cdot 27=19\,683$ entradas; libres $27^{2}\cdot 26=729\cdot 26=18\,954$.
\end{enumerate}
\resultado \textbf{729 entradas (702 libres) para bigramas y 19\,683 (18\,954 libres) para trigramas.}
\comprobacion cada token adicional de contexto multiplica la tabla por $K=27$: $729\cdot 27=19\,683$.
\end{solucion}

\begin{solucion}{eje:2-vf}
\planteamiento contrastar cada afirmación con la definición correspondiente.
\pasos
\begin{enumerate}[label=(\alph*)]
  \item \textbf{Falsa.} El \cref{teo:cadena} es una identidad que se deduce de la regla del producto. La aproximación aparece al imponer la hipótesis de Markov (\cref{def:ngramas}).
  \item \textbf{Verdadera.} Cada fila $i$ es la distribución $P(\cdot\mid i)$, luego $\sum_j P_{ij}=1$. En la \cref{fig:cadena}, la fila \texttt{a} da $3/5+1/5+1/5=1$.
\end{enumerate}
\resultado \textbf{(a) falsa; (b) verdadera.}
\end{solucion}

\begin{solucion}{eje:2-sol-matriz}
\datos corpus «sol», «sal», «los»; $\V=\{\fin,\texttt{a},\texttt{l},\texttt{o},\texttt{s}\}$.
\planteamiento añadir \fin\ al principio y al final de cada palabra, extraer los bigramas, contar y dividir cada fila por su total (\cref{teo:emv}).
\pasos
\begin{enumerate}[label=\arabic*.]
  \item Secuencias: \fin sol\fin, \fin sal\fin, \fin los\fin.
  \item Bigramas (12 en total):
  \[\begin{aligned}
  &\text{sol: } (\fin,s),\,(s,o),\,(o,l),\,(l,\fin);\qquad \text{sal: } (\fin,s),\,(s,a),\,(a,l),\,(l,\fin);\\
  &\text{los: } (\fin,l),\,(l,o),\,(o,s),\,(s,\fin).
  \end{aligned}\]
  \item Matriz de conteos y matriz de transición:
\end{enumerate}
\begin{center}\small
\begin{tabular}{@{}c|ccccc|c@{}}
\toprule
$N$ & \fin & a & l & o & s & total \\ \midrule
\fin & 0 & 0 & 1 & 0 & 2 & 3 \\
a & 0 & 0 & 1 & 0 & 0 & 1 \\
l & 2 & 0 & 0 & 1 & 0 & 3 \\
o & 0 & 0 & 1 & 0 & 1 & 2 \\
s & 1 & 1 & 0 & 1 & 0 & 3 \\ \bottomrule
\end{tabular}
\qquad
\begin{tabular}{@{}c|ccccc@{}}
\toprule
$P$ & \fin & a & l & o & s \\ \midrule
\fin & 0 & 0 & $1/3$ & 0 & $2/3$ \\
a & 0 & 0 & $1$ & 0 & 0 \\
l & $2/3$ & 0 & 0 & $1/3$ & 0 \\
o & 0 & 0 & $1/2$ & 0 & $1/2$ \\
s & $1/3$ & $1/3$ & 0 & $1/3$ & 0 \\ \bottomrule
\end{tabular}
\end{center}
\resultado \textbf{las dos matrices anteriores}; por ejemplo, $P(\texttt{s}\mid\fin)=2/3$ y $P(\texttt{l}\mid\texttt{a})=1$.
\comprobacion los totales suman $3+1+3+2+3=12$ bigramas y cada fila de $P$ suma~1.
\end{solucion}

\begin{solucion}{eje:2-sol-prob}
\datos la matriz $P$ del \cref{eje:2-sol-matriz}.
\planteamiento por el \cref{cor:cadena-secuencias}, la probabilidad de una palabra es el producto de sus transiciones.
\pasos
\begin{enumerate}[label=\arabic*.]
  \item $P(\text{sol})=P(\texttt{s}\mid\fin)P(\texttt{o}\mid\texttt{s})P(\texttt{l}\mid\texttt{o})P(\fin\mid\texttt{l})=\frac23\cdot\frac13\cdot\frac12\cdot\frac23=\frac{4}{54}=\frac{2}{27}$.
  \item $P(\text{sal})=\frac23\cdot\frac13\cdot 1\cdot\frac23=\frac{4}{27}$.
  \item $P(\text{los})=P(\texttt{l}\mid\fin)P(\texttt{o}\mid\texttt{l})P(\texttt{s}\mid\texttt{o})P(\fin\mid\texttt{s})=\frac13\cdot\frac13\cdot\frac12\cdot\frac13=\frac{1}{54}$.
  \item $P(\text{sos})=\frac23\cdot\frac13\cdot\frac12\cdot\frac13=\frac{2}{54}=\frac{1}{27}$.
  \item Cociente: $P(\text{sal})/P(\text{sol})=(4/27)/(2/27)=2$. Las dos palabras solo difieren en el tramo central: $\texttt{s}\to\texttt{a}\to\texttt{l}$ vale $\frac13\cdot 1=\frac13$ y $\texttt{s}\to\texttt{o}\to\texttt{l}$ vale $\frac13\cdot\frac12=\frac16$.
\end{enumerate}
\resultado \textbf{$P(\text{sol})=2/27\approx 0{,}07407$; $P(\text{sal})=4/27\approx 0{,}1481$; $P(\text{los})=1/54\approx 0{,}01852$; $P(\text{sos})=1/27\approx 0{,}03704$.}
\comprobacion «sos» no está en el corpus y tiene probabilidad positiva; la suma de las cuatro es menor que~1, como corresponde a cuatro palabras entre infinitas.
\end{solucion}

\begin{solucion}{eje:2-longitud}
\datos $P(\text{a})=0{,}40$, $P(\text{ana})=0{,}08$, $P(\text{luna})=0{,}10$ (\cref{ej:bigramas}).
\planteamiento cada transición multiplica por un factor $\le 1$; para comparar longitudes distintas se usa la log-probabilidad media por transición, $\frac1n\ln P$.
\pasos
\begin{enumerate}[label=\arabic*.]
  \item «a» tiene 2 transiciones: $\ln 0{,}40/2=-0{,}9163/2=-0{,}4581$.
  \item «ana» tiene 4: $\ln 0{,}08/4=-2{,}526/4=-0{,}6314$.
  \item «luna» tiene 5: $\ln 0{,}10/5=-2{,}303/5=-0{,}4605$.
\end{enumerate}
\resultado \textbf{«a» es la más probable solo por ser la más corta}. Por transición, «a» ($-0{,}4581$) y «luna» ($-0{,}4605$) prácticamente empatan y «ana» queda por detrás. La medida propuesta es la pérdida del \cref{tema:5} cambiada de signo.
\comprobacion en perplejidad por transición, $e^{0{,}4581}\approx 1{,}581$ y $e^{0{,}4605}\approx 1{,}585$.
\end{solucion}

\begin{solucion}{eje:2-sin-fin}
\datos cuatro frases sin símbolo de fin; vocabulario $\{\texttt{a},\texttt{b}\}$ más $\langle s\rangle$.
\planteamiento contar las transiciones desde cada token que tiene sucesor y normalizar; después multiplicar.
\pasos
\begin{enumerate}[label=\arabic*.]
  \item Desde $\langle s\rangle$: a (frases 1 y 4) y b (frases 2 y 3); $P(\texttt{a}\mid\langle s\rangle)=P(\texttt{b}\mid\langle s\rangle)=1/2$.
  \item Desde \texttt{a} (solo cuentan las \texttt{a} seguidas de otro token, frases 1 y 4): $\texttt{a}\to\texttt{b}$ dos veces; $P(\texttt{b}\mid\texttt{a})=1$, $P(\texttt{a}\mid\texttt{a})=0$.
  \item Desde \texttt{b} (frases 2 y 3): $\texttt{b}\to\texttt{b}$ y $\texttt{b}\to\texttt{a}$; $P(\texttt{b}\mid\texttt{b})=P(\texttt{a}\mid\texttt{b})=1/2$.
  \item (b) $P(\texttt{aa})=\frac12\cdot 0=0$; $P(\texttt{ab})=\frac12\cdot 1=\frac12$; $P(\texttt{ba})=\frac12\cdot\frac12=\frac14$; $P(\texttt{bb})=\frac14$. Suma: $0+\frac12+\frac14+\frac14=1$.
  \item (c) Es el caso $L=2$ de la \cref{prop:longitud-fija}: cada longitud suma~1, de modo que la suma sobre todas las longitudes diverge.
\end{enumerate}
\resultado \textbf{$P(\texttt{aa})=0$, $P(\texttt{ab})=1/2$, $P(\texttt{ba})=P(\texttt{bb})=1/4$}; el modelo sin símbolo de fin no es un modelo de lenguaje en el sentido de la \cref{def:ml}.
\comprobacion las ocho frases de longitud~3 también suman~1 (por ejemplo, $P(\texttt{abb})=\frac12\cdot 1\cdot\frac12=\frac14$).
\end{solucion}

% ---------------------------------------------------------------------
\section{Tema 3}

\begin{solucion}{eje:3-laplace}
\datos conteos del \cref{eje:2-sol-matriz}; $\alpha=1$; $K=5$.
\planteamiento $\widehat P_1(v\mid u)=(N(u,v)+1)/(N(u)+5)$ (\cref{def:laplace}).
\pasos
\begin{enumerate}[label=\arabic*.]
  \item $P(\fin\mid\fin)$: $N(\fin,\fin)=0$ y $N(\fin)=3$; $P=(0+1)/(3+5)=1/8$.
  \item $P(\texttt{l}\mid\texttt{a})$: $N(\texttt{a},\texttt{l})=1$ y $N(\texttt{a})=1$; $P=(1+1)/(1+5)=2/6=1/3$.
\end{enumerate}
\resultado \textbf{$P(\fin\mid\fin)=1/8=0{,}125$ (antes~0) y $P(\texttt{l}\mid\texttt{a})=1/3\approx 0{,}3333$ (antes~1).}
\comprobacion la fila \texttt{a} suavizada es $(1,1,2,1,1)/6$, que suma $6/6=1$.
\end{solucion}

\begin{solucion}{eje:3-u1}
\datos cuatro frases; $\alpha=1$.
\planteamiento $P(\text{Ana}\mid\text{soy})=(N(\text{soy},\text{Ana})+1)/(N(\text{soy})+K)$.
\pasos
\begin{enumerate}[label=\arabic*.]
  \item Vocabulario: $\{\langle s\rangle,\text{yo},\text{soy},\text{Ana},\langle/s\rangle,\text{no},\text{quiero},\text{sopa},\text{ni}\}$, luego $K=9$.
  \item «soy» aparece en las frases 1, 2 y 3, siempre con sucesor: Ana (frases 1 y 3) o $\langle/s\rangle$ (frase 2). Así, $N(\text{soy})=3$ y $N(\text{soy},\text{Ana})=2$.
  \item $P=(2+1)/(3+9)=3/12=1/4$.
\end{enumerate}
\resultado \textbf{$P(\text{Ana}\mid\text{soy})=1/4=0{,}25$} (sin suavizar, $2/3\approx 0{,}6667$).
\comprobacion la fila «soy» suavizada asigna $3/12$ a Ana, $2/12$ a $\langle/s\rangle$ y $1/12$ a cada uno de los otros 7 tokens: $3+2+7=12$.
\end{solucion}

\begin{solucion}{eje:3-interp}
\datos $N(\text{soy},\text{Ana})=2$, $N(\text{soy})=3$, $N(\text{Ana})=4$; total de tokens $5+5+5+8=23$; $\lambda_1=\lambda_2=1/2$.
\planteamiento $P_{\mathrm{int}}=\tfrac12\widehat P_2(\text{Ana}\mid\text{soy})+\tfrac12\widehat P_1(\text{Ana})$.
\pasos
\begin{enumerate}[label=\arabic*.]
  \item Bigrama: $\widehat P_2=2/3$.
  \item Unigrama: «Ana» aparece una vez en cada frase, $\widehat P_1=4/23$.
  \item $P_{\mathrm{int}}=\tfrac12\cdot\tfrac23+\tfrac12\cdot\tfrac{4}{23}=\tfrac13+\tfrac{2}{23}=\tfrac{23}{69}+\tfrac{6}{69}=\tfrac{29}{69}$.
\end{enumerate}
\resultado \textbf{$P(\text{Ana}\mid\text{soy})=29/69\approx 0{,}4203$.}
\comprobacion el valor está entre el unigrama ($0{,}1739$) y el bigrama ($0{,}6667$), como corresponde a una media ponderada.
\end{solucion}

\begin{solucion}{eje:3-limite}
\datos $u$ fijo con $n=N(u)>0$; $\widehat P_\alpha(v\mid u)=(N(u,v)+\alpha)/(n+\alpha K)$.
\planteamiento límites de un cociente de funciones continuas en $\alpha$.
\pasos
\begin{enumerate}[label=\arabic*.]
  \item $\alpha\to 0^{+}$: numerador $\to N(u,v)$ y denominador $\to n>0$; el cociente tiende a $N(u,v)/n$, que es el EMV del \cref{teo:emv}.
  \item $\alpha\to\infty$: dividiendo numerador y denominador por $\alpha$, $\widehat P_\alpha=(N(u,v)/\alpha+1)/(n/\alpha+K)\to 1/K$.
  \item Interpretación: $\widehat P_\alpha$ es la frecuencia relativa que se obtendría si, además de los datos reales, se hubiera observado $\alpha$ veces cada uno de los $K$ tokens tras $u$. Es la media a posteriori con prior de Dirichlet (\cref{prop:laplace}).
\end{enumerate}
\resultado \textbf{$\lim_{\alpha\to 0^+}\widehat P_\alpha=N(u,v)/N(u)$ y $\lim_{\alpha\to\infty}\widehat P_\alpha=1/K$}. Un $\alpha$ grande da más peso a la creencia previa uniforme que a los datos.
\comprobacion con la fila \texttt{a} del \cref{eje:3-laplace} y $\alpha=100$: $P(\texttt{l}\mid\texttt{a})=101/501\approx 0{,}2016$, cerca de $1/5$.
\end{solucion}

% ---------------------------------------------------------------------
\section{Tema 4}

\begin{solucion}{eje:4-orden}
\planteamiento aplicar la propiedad (iii) de la \cref{prop:softmax}.
\pasos
\begin{enumerate}[label=\arabic*.]
  \item Para todo $T>0$, $s\mapsto e^{s/T}$ es estrictamente creciente.
  \item Si $z_i>z_j$, entonces $e^{z_i/T}>e^{z_j/T}$ y, al dividir por el mismo denominador, $p_i>p_j$.
\end{enumerate}
\resultado \textbf{No.} La temperatura modifica cuánto destaca el máximo, no cuál es.
\comprobacion en el \cref{tab:temperatura}, $p_1$ es la mayor con $T=0{,}5$, $1$ y $2$.
\end{solucion}

\begin{solucion}{eje:4-softmax}
\datos $z=(0,\ln 2,\ln 5)$; $T=1$ y $T=2$.
\planteamiento \cref{def:softmax}, usando $e^{\ln a}=a$ y $e^{(\ln a)/2}=\sqrt a$.
\pasos
\begin{enumerate}[label=\arabic*.]
  \item $T=1$: $e^{z}=(1,2,5)$, suma $8$; $\softmax(z)=(1/8,\,2/8,\,5/8)=(0{,}125;\,0{,}25;\,0{,}625)$.
  \item $T=2$: $e^{z/2}=(1,\sqrt2,\sqrt5)=(1;\,1{,}414;\,2{,}236)$, suma $4{,}650$.
  \item Cocientes: $1/4{,}650=0{,}2151$; $1{,}414/4{,}650=0{,}3041$; $2{,}236/4{,}650=0{,}4809$.
  \item $T\to 0^{+}$: el máximo $z_3=\ln 5$ es único; por la propiedad (iv), el límite es $(0,0,1)$.
\end{enumerate}
\resultado \textbf{$T=1$: $(0{,}125;\,0{,}25;\,0{,}625)$. $T=2$: $(0{,}2151;\,0{,}3041;\,0{,}4809)$. $T\to 0^{+}$: $(0,0,1)$.}
\comprobacion la segunda fila suma $0{,}2151+0{,}3041+0{,}4809=1{,}000$; al pasar de $T=1$ a $T=2$ el máximo baja de $0{,}625$ a $0{,}4809$.
\end{solucion}

\begin{solucion}{eje:4-muestreo}
\datos $p=(0{,}1;\,0{,}5;\,0{,}4)$; $u\in\{0{,}05;\,0{,}35;\,0{,}6;\,0{,}95\}$.
\planteamiento $X=\min\{k:u<F_k\}$ (\cref{prop:muestreo}).
\pasos
\begin{enumerate}[label=\arabic*.]
  \item $F_1=0{,}1$; $F_2=0{,}6$; $F_3=1{,}0$. Intervalos: $x\leftrightarrow[0;0{,}1)$, $y\leftrightarrow[0{,}1;0{,}6)$, $z\leftrightarrow[0{,}6;1)$.
  \item $u=0{,}05<0{,}1$: \textbf{x}. $u=0{,}35$: $0{,}1\le 0{,}35<0{,}6$, \textbf{y}.
  \item $u=0{,}6$: $0{,}6$ no es menor que $F_2=0{,}6$, pero sí que $F_3$: \textbf{z}. $u=0{,}95$: \textbf{z}.
\end{enumerate}
\resultado \textbf{x, y, z, z.}
\comprobacion la longitud de cada intervalo coincide con su probabilidad: $0{,}1$, $0{,}5$ y $0{,}4$.
\end{solucion}

\begin{solucion}{eje:4-potencia}
\datos $Q(i)\propto P(i)^{\beta}$; $P=(0{,}5;\,0{,}25;\,0{,}25)$.
\planteamiento escribir $P(i)^{\beta}=e^{\beta\ln P(i)}=e^{z_i/T}$ con $z_i=\ln P(i)$ y $T=1/\beta$.
\pasos
\begin{enumerate}[label=\arabic*.]
  \item (a) Duplicar la temperatura ($T=2$) es $\beta=1/2$; reducirla a la mitad ($T=1/2$) es $\beta=2$.
  \item (b) $\beta=2$: $P^{2}=(0{,}25;\,0{,}0625;\,0{,}0625)$, suma $0{,}375$; $Q=(2/3;\,1/6;\,1/6)$.
  \item $\beta=1/2$: $\sqrt{P}=(0{,}7071;\,0{,}5;\,0{,}5)$, suma $1{,}707$; $Q=(0{,}4142;\,0{,}2929;\,0{,}2929)$.
\end{enumerate}
\resultado \textbf{$\beta=1/T$}; con $\beta=2$, $Q=(0{,}6667;\,0{,}1667;\,0{,}1667)$; con $\beta=1/2$, $Q=(0{,}4142;\,0{,}2929;\,0{,}2929)$.
\comprobacion ambas distribuciones suman~1 y conservan el orden de $P$.
\end{solucion}

\begin{solucion}{eje:4-limites}
\datos $z\in\R^{K}$, $T>0$, $c\in\R$.
\planteamiento manipular el cociente y tomar límites término a término.
\pasos
\begin{enumerate}[label=\arabic*.]
  \item (ii) $e^{(z_i+c)/T}=e^{z_i/T}e^{c/T}$ por $e^{a+b}=e^{a}e^{b}$. El factor $e^{c/T}>0$ aparece en el numerador y en cada sumando del denominador; al sacarlo como factor común se cancela.
  \item (iv) Sea $z_k$ el máximo único. Con $c=-z_k$: $p_i=e^{(z_i-z_k)/T}/\sum_j e^{(z_j-z_k)/T}$. Para $j=k$ el sumando vale~1; para $j\ne k$, $z_j-z_k<0$, luego $(z_j-z_k)/T\to-\infty$ y $e^{(z_j-z_k)/T}\to 0$. El denominador tiende a~1, $p_k\to 1$ y $p_i\to 0$ si $i\ne k$.
  \item (v) Para $z$ fijo, $z_j/T\to 0$ y, por continuidad, $e^{z_j/T}\to 1$; luego $p_i\to 1/K$.
  \item Contraejemplo para (iv): $z=(1,1,0)$. Con $c=-1$, $p=(1,1,e^{-1/T})/(2+e^{-1/T})\to(1/2,1/2,0)$, que no es un vector canónico.
\end{enumerate}
\resultado \textbf{Las tres propiedades quedan demostradas; sin máximo único, el límite reparte la masa entre los máximos.} $\blacksquare$
\comprobacion con $z=(2,1,0)$ y $T=0{,}1$, $\softmax_T(z)\approx(0{,}99995;\,4{,}54\cdot 10^{-5};\,2\cdot 10^{-9})$; con $T=100$, $\approx(0{,}3367;\,0{,}3333;\,0{,}3300)$.
\end{solucion}

\begin{solucion}{eje:4-lse}
\datos $g(x)=\ln\sum_i e^{x_i}$; $x=(1000;\,1001;\,1002)$.
\planteamiento regla de la cadena para (a); factor común para (b); (b) con $b=-\max_i x_i$ para (c).
\pasos
\begin{enumerate}[label=\arabic*.]
  \item (a) Con $\frac{d}{du}\ln u=\frac1u$: $\partial g/\partial x_i=e^{x_i}/\sum_j e^{x_j}=\softmax(x)_i$.
  \item (b) $\sum_i e^{x_i+b}=e^{b}\sum_i e^{x_i}$; tomando logaritmos, $g(x+b\mathbf 1)=b+g(x)$.
  \item (c) Con $m=1002$: $g(x)=1002+\ln(e^{-2}+e^{-1}+e^{0})$.
  \item $e^{-2}=0{,}1353$; $e^{-1}=0{,}3679$; suma $0{,}1353+0{,}3679+1=1{,}503$; $\ln 1{,}503=0{,}4076$.
  \item $g(x)=1002+0{,}4076=1002{,}408$.
  \item Directamente: $e^{1000}\approx 10^{434}$ supera el máximo de doble precisión ($\approx 1{,}8\cdot 10^{308}$); se obtiene $+\infty$ y $g=+\infty$, un resultado falso.
\end{enumerate}
\resultado \textbf{$g(x)\approx 1002{,}408$}; el cálculo directo desborda. Además, se obtiene la softmax sin coste adicional:
\[\softmax(x)=\frac{(0{,}1353;\,0{,}3679;\,1)}{1{,}503}=(0{,}09003;\,0{,}2447;\,0{,}6652).\]
\comprobacion por la \cref{prop:lse}, $1002\le 1002{,}408\le 1002+\ln 3=1003{,}099$.
\end{solucion}

% ---------------------------------------------------------------------
\section{Tema 5}

\begin{solucion}{eje:5-ppl1}
\planteamiento propiedad (ii) de la \cref{prop:ppl}.
\pasos
\begin{enumerate}[label=\arabic*.]
  \item $\PPL=1$ equivale a $\mathcal J=0$.
  \item $\mathcal J$ es la media de los términos $-\ln p_t\ge 0$; es nula si y solo si todos lo son, es decir, $p_t=1$ para todo $t$.
\end{enumerate}
\resultado \textbf{El modelo predijo cada token del texto con probabilidad~1.}
\end{solucion}

\begin{solucion}{eje:5-perdida}
\datos $p=(1/2,\,1/4,\,1/4,\,1/8)$; $N=4$.
\planteamiento en base~2 los logaritmos son enteros; después se convierte a nats multiplicando por $\ln 2$.
\pasos
\begin{enumerate}[label=\arabic*.]
  \item $-\log_2 p_t=1,\,2,\,2,\,3$; suma $8$; media $\mathcal J=2$ bits.
  \item En nats: $2\ln 2=2\cdot 0{,}6931=1{,}386$.
  \item $\PPL=2^{2}=4$; equivalentemente, $e^{1{,}386}=4{,}000$.
\end{enumerate}
\resultado \textbf{$\mathcal J=2$ bits $=1{,}386$ nats; $\PPL=4$.}
\comprobacion por la \cref{prop:ppl}(i), $\PPL=(2\cdot 4\cdot 4\cdot 8)^{1/4}=256^{1/4}=4$.
\end{solucion}

\begin{solucion}{eje:5-digitos}
\datos $P(0)=0{,}82$; $P(d)=0{,}02$ para $d=1,\dots,9$; prueba con nueve ceros y un~3, $N=10$.
\planteamiento $\PPL=\exp\bigl(-\frac1N\sum_t\ln p_t\bigr)$.
\pasos
\begin{enumerate}[label=\arabic*.]
  \item $\ln 0{,}82=-0{,}1985$; para nueve ceros, $9\cdot(-0{,}1985)=-1{,}786$.
  \item $\ln 0{,}02=-3{,}912$.
  \item Suma $-1{,}786-3{,}912=-5{,}698$; pérdida media $5{,}698/10=0{,}5698$ nats.
  \item $\PPL=e^{0{,}5698}=1{,}768$. El modelo uniforme tiene $\PPL=10$ por la \cref{prop:ppl}(iii).
\end{enumerate}
\resultado \textbf{$\PPL\approx 1{,}768$ frente a~10 del modelo uniforme.}
\comprobacion $(0{,}82^{9}\cdot 0{,}02)^{-1/10}=(0{,}1676\cdot 0{,}02)^{-0{,}1}=0{,}003352^{-0{,}1}\approx 1{,}768$.
\end{solucion}

\begin{solucion}{eje:5-ppl}
\datos $p_1,\dots,p_N\in(0,1]$.
\planteamiento propiedades del logaritmo: $\ln(1/x)=-\ln x$, $\sum\ln a_t=\ln\prod a_t$, $a\ln x=\ln x^{a}$, $e^{\ln y}=y$.
\pasos
\begin{enumerate}[label=\arabic*.]
  \item $\mathcal J=\frac1N\sum_t\ln(1/p_t)=\frac1N\ln\prod_t(1/p_t)=\ln\bigl(\prod_t 1/p_t\bigr)^{1/N}$.
  \item $\PPL=e^{\mathcal J}=\bigl(\prod_t 1/p_t\bigr)^{1/N}$.
  \item Como $0<p_t\le 1$, cada $1/p_t\ge 1$; el producto es $\ge 1$ y, al ser $x\mapsto x^{1/N}$ creciente con $1^{1/N}=1$, también su raíz. La igualdad exige que todos los factores valgan~1.
  \item Si $p_t=1/K$ para todo $t$: $\PPL=(K^{N})^{1/N}=K$.
\end{enumerate}
\resultado \textbf{Las tres afirmaciones de la \cref{prop:ppl}.} $\blacksquare$
\comprobacion \cref{eje:5-perdida}: $(2\cdot 4\cdot 4\cdot 8)^{1/4}=4=2^{2}$.
\end{solucion}

\begin{solucion}{eje:5-lagrange}
\datos $p$ fija con $p(x)>0$ y $\sum_x p(x)=1$; se busca $q$ con $\sum_x q(x)=1$ que minimice $H(p,q)$.
\planteamiento minimización con una restricción de igualdad.
\pasos
\begin{enumerate}[label=\arabic*.]
  \item Lagrangiana: $\Lambda(q,\lambda)=-\sum_x p(x)\ln q(x)+\lambda\bigl(\sum_x q(x)-1\bigr)$.
  \item $\partial\Lambda/\partial q(x)=-p(x)/q(x)+\lambda=0$, luego $q(x)=p(x)/\lambda$.
  \item Restricción: $\sum_x p(x)/\lambda=1$, luego $\lambda=1$ y $q=p$.
  \item Es un mínimo global: $\partial^{2}H/\partial q(x)^{2}=p(x)/q(x)^{2}>0$ y las derivadas cruzadas son nulas, de modo que $H(p,\cdot)$ es estrictamente convexa sobre el conjunto convexo de distribuciones.
  \item Valor mínimo: $H(p,p)=H(p)$.
\end{enumerate}
\resultado \textbf{$H(p,q)$ alcanza su mínimo, $H(p)$, únicamente en $q=p$.} $\blacksquare$
\comprobacion con $p=(0{,}5;\,0{,}3;\,0{,}2)$: $H(p)=0{,}3466+0{,}3612+0{,}3219=1{,}030$ nats; con $q=(0{,}6;\,0{,}2;\,0{,}2)$ se obtiene $1{,}060$, y con la uniforme, $\ln 3=1{,}099$.
\end{solucion}


\backmatter
% =====================================================================
\begin{thebibliography}{99}
\addcontentsline{toc}{chapter}{Bibliografía}
\markboth{Bibliografía}{}
\small

\bibitem{bt}
D.~P. Bertsekas y J.~N. Tsitsiklis.
\newblock \emph{Introduction to Probability}, 2.ª ed.
\newblock Athena Scientific, 2008.
\newblock Texto del curso 6.041 del MIT.

\bibitem{bishop}
C.~M. Bishop.
\newblock \emph{Pattern Recognition and Machine Learning}.
\newblock Springer, 2006.

\bibitem{chen-goodman}
S.~F. Chen y J.~Goodman.
\newblock An empirical study of smoothing techniques for language modeling.
\newblock \emph{Computer Speech \& Language}, 13(4):359--394, 1999.

\bibitem{ct}
T.~M. Cover y J.~A. Thomas.
\newblock \emph{Elements of Information Theory}, 2.ª ed.
\newblock Wiley, 2006.

\bibitem{cs224n}
D.~Yang y Y.~Choi.
\newblock CS224N: Natural Language Processing with Deep Learning.
\newblock Stanford University, invierno de 2026. Clase~4, «Language Models and RNNs».
\newblock \url{https://web.stanford.edu/class/cs224n/}

\bibitem{fan}
A.~Fan, M.~Lewis e Y.~Dauphin.
\newblock Hierarchical neural story generation.
\newblock En \emph{Proceedings of ACL}, 2018.
\newblock \url{https://arxiv.org/abs/1805.04833}

\bibitem{gbc}
I.~Goodfellow, Y.~Bengio y A.~Courville.
\newblock \emph{Deep Learning}.
\newblock MIT Press, 2016.
\newblock \url{https://www.deeplearningbook.org/}

\bibitem{holtzman}
A.~Holtzman, J.~Buys, L.~Du, M.~Forbes e Y.~Choi.
\newblock The curious case of neural text degeneration.
\newblock En \emph{International Conference on Learning Representations (ICLR)}, 2020.
\newblock \url{https://arxiv.org/abs/1904.09751}

\bibitem{jm3}
D.~Jurafsky y J.~H. Martin.
\newblock \emph{Speech and Language Processing}, 3.ª ed., borrador del 19 de agosto de 2026.
\newblock Cap.~3, «N-gram Language Models».
\newblock \url{https://web.stanford.edu/~jurafsky/slp3/}

\bibitem{mackay}
D.~J.~C. MacKay.
\newblock \emph{Information Theory, Inference, and Learning Algorithms}.
\newblock Cambridge University Press, 2003.
\newblock \url{http://www.inference.org.uk/itila/book.html}

\bibitem{mit6041}
J.~N. Tsitsiklis.
\newblock 6.041SC: Probabilistic Systems Analysis and Applied Probability.
\newblock Massachusetts Institute of Technology, MIT OpenCourseWare, otoño de 2013.
\newblock \url{https://ocw.mit.edu/courses/6-041sc-probabilistic-systems-analysis-and-applied-probability-fall-2013/}

\bibitem{shannon}
C.~E. Shannon.
\newblock A mathematical theory of communication.
\newblock \emph{Bell System Technical Journal}, 27:379--423 y 623--656, 1948.

\bibitem{d2l}
A.~Zhang, Z.~C. Lipton, M.~Li y A.~J. Smola.
\newblock \emph{Dive into Deep Learning}, versión 1.0.3.
\newblock Cambridge University Press, 2023.
\newblock \url{https://d2l.ai/}

\end{thebibliography}


\end{document}
